\section{Local BHSM topography on the hyperspherical parent}
\label{sec:bhsm-los-topography}

The smooth $S^3(R_H)$ background and the physical $n=2$ response do not exhaust the geometry sampled by a photon. We therefore separate the global parent from a local deformation field $\psi(\mathbf x,t)$:
\begin{equation}
\gamma_{ij}=a^2 e^{2\psi}\bar\gamma^{S^3}_{ij}.
\end{equation}
To first order its intrinsic-curvature perturbation is
\begin{equation}
\delta{}^{(3)}R=-\frac{4}{a^2}\left(\bar\nabla^2+\frac{3}{R_H^2}\right)\psi,
\end{equation}
so the physical manuscript $n=2$ harmonic carries the shifted factor $8-3=5$.  The local volume rate is
\begin{equation}
H_{\rm local}=H+\dot\psi.
\end{equation}

The retained BHSM topographic flux candidate is $F_T=\nabla T-B\nabla(\nabla^2T)$ with $\nabla\cdot F_T=S$.  The existing seam-charge candidate is
\begin{equation}
V_\Sigma^2=\frac{c^2}{8\pi}\int_{\Sigma_\star}H_\Sigma\Delta K_{uu}\,dA,
\end{equation}
giving $V_\Sigma^2=GM/r_\star$ at leading spherical weak-field order.  The exact action-level export from this boundary charge into the dynamical $\psi$ source remains open; no normalization is invented here.

\subsection{Localized line-of-sight response}
For a minimal coherent local-source proxy the accumulated radial perturbation grows only to a coherence distance $\chi_c$ and then saturates.  The closed-$S^3$ radial distance term therefore has the normalized shape
\begin{equation}
\mathcal F_{\rm loc}(z)=\frac{[C_K/S_K](\chi)\min(\chi,\chi_c)}{[C_K/S_K](\chi_\star)\min(\chi_\star,\chi_c)}.
\end{equation}
The value $z_c=0.03$ originated as an illustrative local-coherence choice
and is not claimed to be action-derived or newly fitted here. It is carried
forward unchanged and is therefore frozen for the retrospective
Section~\ref{sec:sn-sightlines} interpretation test.
The table below retains the historical $-0.0412$ mag normalization;
it is not the later full-covariance calibration.

\begin{center}
\begin{tabular}{r r}
$z$ & $A_\mu^{\rm loc}$ [mag]\\
\hline
0.02 & -0.041200\\
0.10 & -0.012598\\
0.30 & -0.004435\\
0.50 & -0.002810\\
0.80 & -0.001904\\
1.00 & -0.001605\\
1.50 & -0.001212\\
2.10 & -0.000991\\
\end{tabular}
\end{center}

\paragraph{Calibration boundary.}
The carried-forward low-redshift SN scalar amplitude normalizes the local topographic response in this exploratory decomposition.  It cannot simultaneously determine an independent global $n=2$ amplitude.  The global/local split must ultimately be fixed by the common BHSM action or an independent preregistered observable.
